\documentclass[../main.tex]{subfiles}

\begin{document}
We are interested in continuing the research of finding globally optimal solutions over compact spaces in both game theory and in robotics.

\subsection*{Inverse Kinematics}
Our approach to reformulating the polynomial formulation in \cref{sec:kinematicsformulation} was very direct. While we have experimented briefly with solvers capable of polynomial optimization by different techniques, such as \emph{RAPOSa}, \emph{Baron} and \emph{Couenne}, this has not resulted in runtime improvements. Still, the structure of the problem should enable us to create constraints which could significantly speed up the solution process. For example, the feasible set of the multilinear kinematic constraints should be a polytope. While its full representation may be too costly to construct, finding a few of its facets near the optimum could be feasible. Similarly, the results of \textcite{trutman2022globally} indicate that the constraints have a lower-degree representation even without auxiliary variables. Whether this extends to higher degrees of freedom is still an open question.

Our method currently seems most appropriate for the use-cases of manipulator design, where determining the actual performance of manipulators in some metric specified by the user is likely to be preferable to a heuristic value. Moreover, the runtime requirements are not so strict in such contexts. Although our method does not apply to real-time control, additional pre-processing and appropriate heuristics may make such a use case possible. There already exist appropriate methods in the literature. For example, we could use modern real-time IK solvers based on convex QP as runtime heuristics or generate initial feasible solutions to warm-start the search. While running the experiments, we observed that a significant fraction of the solution time for \emph{Gurobi} was spent proving that the solution was optimal. Using the outer approximation technique described by \textcite{Dai2019-uf} could make this stage faster. Finally, it could be possible to use the current configuration of a manipulator to construct a feasible solution using homotopy methods. To make use of such homotopies would likely involve reformulating the kinematic constraint using quaternions.

In practice, we are also able to significantly accelerate the proof-of-concept technique of \textcite{trutman2022globally}. Their Grobner basis computation is implemented directly in \emph{Maple}. Including a cheap pre-processing step to find the required number of polynomials (and moving to \emph{Msolve}) is a simple way to optimize the slowest step of the algorithm.

\subsection*{Game Theory}
The convergence of our algorithm for separable network games was experimentally verified on selected examples. However, to prove that the hierarchy in fact converges is an interesting open problem for further research. While running the experiments, we noticed the computed payoffs were generally of high quality regardless of the solver used; The corresponding measures, on the other hand, were less accurate. For example, the strategies recovered in \cref{ex:separable.security} yields a loss of only $0.16$ instead of $0.17$ for one of the players. This appears to be a problem inherent to our approach rather than something that could be fine-tuned within the solvers.

The idea of the multiple oracle algorithm is to construct a sequence of finite subgames whose equilibria approximate the equilibrium of a given continuous game in the Wasserstein metric. Note that the multiple oracle algorithm makes it possible to approximate the equilibrium of \emph{any} continuous games in the sense of Theorems \ref{thm:convergence1} and \ref{thm:convergence2}, with the caveat that an individual sequence of equilibria may fail to converge. Although possible in theory \cite{Adam21}, this behavior has never been observed in the sample games.
\end{document}